\documentclass[11pt,a4paper]{article}
\usepackage[utf8]{inputenc}
\usepackage[T1]{fontenc}
\usepackage[portuguese]{babel}
\usepackage{xcolor}
\usepackage{hyperref}
\usepackage{geometry}
\usepackage{tcolorbox}

\geometry{margin=2.5cm}

\definecolor{primary}{RGB}{41, 128, 185}
\definecolor{secondary}{RGB}{44, 62, 80}
\definecolor{codebg}{RGB}{240, 240, 240}

\hypersetup{colorlinks=true, linkcolor=primary, urlcolor=primary}

\title{\color{secondary}\Huge\textbf{Exercício 4.6: O projeto v2.0}}
\author{\color{primary}DevOps com Kubernetes -- 2026}
\date{}

\begin{document}
\maketitle

\section*{\color{primary}Instruções do Exercício}
\begin{tcolorbox}[colback=blue!5!white,colframe=blue!75!black,title=Exercício 4.6: O projeto v2.0]
\textit{Nota: Este exercício aborda a escalabilidade do broadcaster de mensagens.}

O \textit{broadcaster} de mensagens (usando NATS) deve ser capaz de ser escalado sem enviar a mensagem várias vezes ao cliente. 

Teste a sua infraestrutura certificando-se de que o sistema consegue rodar com 6 réplicas ativas do broadcaster sem problemas ou travamentos. 

Lembre-se de que as mensagens só precisam ser enviadas para o serviço externo se todos os serviços estiverem funcionando corretamente. Uma mensagem perdida aleatoriamente não é um problema grave, mas o envio de uma duplicata é.
\end{tcolorbox}

\end{document}
